\section{FFN/MoE vs. Attention：计算角色的本质差异}

\subsection{两种计算范式}

\begin{importantbox}{Attention vs. FFN 的本质区别}
\begin{itemize}
  \item \textbf{Attention}：刻画 token \textbf{之间}的交互（interaction），是因果的（causal mask），当前 query 只能看到过去的 key
  \item \textbf{FFN/MoE}：是 \textbf{token 级别}的映射（projection），token 之间没有交互，每个 token 独立地经过 FFN 变换
\end{itemize}
\end{importantbox}

这个区别有重要的工程意义：因为 MoE/FFN 是 token-level 的操作，所以：
\begin{itemize}
  \item 路由是\textbf{per-token} 的：每个 token 独立选择 top-K Expert
  \item 可以按 Expert 对 token 进行分桶（grouping），将路由到同一 Expert 的 token 聚合后批量计算
  \item 不涉及 token 间的依赖关系，天然适合并行
\end{itemize}

\subsection{参数量占比分析}

\begin{table}[H]
\centering
\caption{GQA vs. MoE/FFN 参数量占比}
\begin{tabular}{lccc}
\toprule
\textbf{模型} & \textbf{GQA 占比} & \textbf{FFN/MoE 占比} & \textbf{说明} \\
\midrule
Qwen3-30B-A3B（全量 128 Expert） & 2.9\% & 94\% & 总参数 30.5B \\
Qwen3-30B-A3B（激活 8 Expert） & 27\% & 54\% & 激活参数 3.3B \\
Qwen3-32B（Dense） & 18\% & 76\% & 总参数 32B \\
\bottomrule
\end{tabular}
\end{table}

\begin{knowledgebox}{FFN/MoE 始终是参数量大户}
无论是 Dense 还是 MoE，FFN 部分始终占据参数量的绝大多数。MoE 的全量参数中，128 个 Expert 贡献了 94\% 的参数量。即使只看激活参数（8 个 Expert），FFN 仍占 54\%。
\end{knowledgebox}

\subsection{本章小结}

Attention 负责 token 间的信息交互，FFN/MoE 负责 token 内部的特征变换。两者功能互补，但参数量分布极不均衡——FFN/MoE 是绝对的大户。

